% Part: first-order-logic
% Chapter: completeness
% Section: downward-ls

\documentclass[../../../include/open-logic-section]{subfiles}

\begin{document}

\olfileid{fol}{com}{dls}
\olsection{لووینہائم--سکولم دا قضیہ}

لووینہائم--سکولم دا قضیہ کہندا اے کہ جے کسے نظریے
دا لامتناہی نمونہ اے، تاں اوہدا اک ایہو جیہا نمونہ
وی اے جیہڑا زیادہ توں زیادہ !!{denumerable} اے۔
ایس گل دا فوری نتیجہ ایہ اے کہ پہلے درجے دی منطق
ایہ بیان نہیں کر سکدی کہ کسے !!a{structure}
دا سائز !!{nonenumerable} اے: کوئی وی
!!{sentence} یا !!{sentence}s دا سیٹ جیہڑا
ساریاں !!{nonenumerable} !!{structure}s وچ پورا
ہوندا اے، کسے !!{enumerable} ساخت وچ وی پورا
ہوندا اے۔

\begin{thm}
\ollabel{thm:downward-ls}
جے $\Gamma$ بے تضاد اے، تاں اوہدا
!!a{enumerable} نمونہ اے؛ یعنی اوہ اک ایسی
!!a{structure} وچ پورا ہون جوگا اے جیہدا
دائرہ یا محدود اے یا !!{denumerable}۔
\end{thm}

\begin{proof}
جے $\Gamma$ بے تضاد اے، تاں تمامیت دے قضیے
دے ثبوت توں ملی !!{structure}~$\Struct M$
% PNB-COM-020: تعمیر Henkin دی ودھی زبان دے بند ٹرماں یا اوہناں دے طبقیاں اُتے اے۔
دا دائرہ $\Domain{M}$ ودھی زبان~$\Lang L'$
دے بند ٹرماں دے سیٹ توں وڈا نہیں؛طبقیاں دی گنتی وی نمائندیاں
دی گنتی توں وڈّی نہیں۔ اصل شمار جوگ علامات وچ شمار جوگ نویاں
مستقل علامتاں شامل کرن نال ایہ بند ٹرماں دا سیٹ شمار جوگ ہی
رہندا اے۔ ایس لئی
$\Struct M$ زیادہ توں زیادہ !!{denumerable}
اے۔
\end{proof}

\begin{thm}
\ollabel{noidentity-ls}
جے $\Gamma$ برابری دے بغیر پہلے درجے دی
منطق دی زبان وچ !!{sentence}s دا اک بے تضاد
سیٹ اے، تاں اوہدا !!a{denumerable} نمونہ اے؛
یعنی اوہ اک ایسی !!a{structure} وچ پورا ہون
جوگا اے جیہدا دائرہ لامتناہی تے
!!{enumerable} اے۔
\end{thm}

\begin{proof}
جے $\Gamma$ بے تضاد اے تے اوہدے وچ کوئی
ایہو جیہا جملہ نہیں جیہدے وچ برابری آوے،
تاں تمامیت دے قضیے دے ثبوت توں ملی
!!{structure}~$\Struct M$ دا دائرہ
% PNB-COM-011: ٹرم نمونے دا اصل دائرہ کھلے ٹرماں دی بجائے بند ٹرم نیں۔
$\Domain{M}$ زبان~$\Lang L'$ دے بند
ٹرماں دے سیٹ دے عین برابر اے۔ سو
$\Struct{M}$ !!{denumerable} اے، کیونکہ
$\Trm[L']$ گنتی وار فہرست بنن جوگ اے،
تے اوس پھیلی زبان وچ شامل لامتناہی نویاں مستقل
علامتاں بند ٹرماں دے سیٹ نوں لامتناہی
بناندیاں نیں۔
\end{proof}

\begin{ex}[سکولم دا پیراڈوکس]
زرمیلو--فرینکل سیٹ تھیوری~$\Log{ZFC}$
اک بہت طاقت ور ڈھانچہ اے جیہدے وچ
تقریباً سارے ریاضیاتی بیان کیتے جا سکدے
نیں، سیٹاں دے سائز بارے حقیقتاں وی۔
مثلاً $\Log{ZFC}$ ایہ ثابت کر سکدا اے
کہ حقیقی اعداد دا سیٹ~$\Real$
!!{nonenumerable} اے؛ ایہ کانتور دا
قضیہ وی ثابت کر سکدا اے کہ کسے وی
سیٹ دے ذیلی سیٹاں دا سیٹ اوس سیٹ
توں وڈا ہوندا اے، وغیرہ۔ جے
$\Log{ZFC}$ بے تضاد اے، تاں اوہدے
سارے نمونے لامتناہی نیں۔ ہور ایہ کہ
اوہناں ساریاں وچ ایہوجے !!{element}s
وی نیں جنہاں بارے نظریہ کہندا اے کہ
اوہ !!{nonenumerable} نیں؛ مثال لئی
اوہ عنصر جیہڑا قدرتی اعداد دے ذیلی
سیٹاں دے سیٹ دے وجود بارے
~$\Log{ZFC}$ دے قضیے نوں سچا کردا
اے۔ لووینہائم--سکولم دے قضیے مطابق
$\Log{ZFC}$ دے !!{enumerable}
نمونے وی نیں---ایہوجے نمونے جنہاں
وچ ”!!{nonenumerable}“ سیٹ ہوندے
نیں، حالانکہ نمونے آپ !!{enumerable}
ہوندے نیں۔
\end{ex}

\end{document}
